\scanpage{184}
\mspage{133}
\[
\begin{aligned}
\text{a)}\quad &\unsure{p|\rho}_0\in W\\
\text{b)}\quad &p_1\text{ et }p_2\in W\\
\text{c)}\quad &p\in W
\end{aligned}
\]
Alors deux de ces conditions impliquent la troisième.

Résulte aussitôt de la \unsure{définition 3|définition 1}, et du résultat similaire (cor. page 126) pour le diagramme cubique similaire, avec $Y,Y'$ remplacés par
\diagram{\begin{tikzcd}[ampersand replacement=\&,row sep=tiny,column sep=small]
\&X_1\\
\mathscr X=\int X_0\arrow[ru]\arrow[rd]\&\\
\&X_2
\end{tikzcd}}
et
\diagram{\begin{tikzcd}[ampersand replacement=\&,row sep=tiny,column sep=small]
\&X'_1\\
\mathscr X'=\unsure{\cdot}\int X'_0\arrow[ru]\arrow[rd]\&\\
\&X'_2 .
\end{tikzcd}}

\unsure{3.3.} \uline{Intégration propre versus intégration lisse}. \uline{Énoncé d'équivalence entre les deux pts de vue}.
% slanted
\margin{$+$ \uline{commentaires sur le \deleted{\unsure{cas}} des limites $W$-homotopiques quelconques}.}

La définition 3, et les corollaires précédents, n'ont d'intérêt que dans la mesure où on dispose de critères maniables, permettant d'affirmer qu'un carré commutatif est bel et bien homotopiquement cocartésien. C'est ce qu'on va essayer de faire tantôt. Il y a aussi une perplexité, concernant le rôle privilégié joué, dans la théorie développée ici, par les immersions \deleted{\unsure{ouvertes}}\unsure{fermées}, à l'exclusion des immersions \deleted{fermées}\inserted{ouvertes}. Nous avons en somme \deleted{\unsure{recon}} reconnu, des
% archive sheet number: 184
